\documentclass[UTF8,11pt]{ctexart}
\usepackage[a4paper,margin=24mm]{geometry}
\usepackage{amsmath,amssymb,amsthm,mathtools,mathrsfs}
\usepackage[all]{xy}
\usepackage{xurl,hyperref,enumitem}
\usepackage{tikz}
\usetikzlibrary{arrows.meta,cd,decorations.markings,calc,patterns}
\hypersetup{hidelinks,breaklinks=true}
\setlength{\emergencystretch}{3em}
\allowdisplaybreaks
\newtheorem*{theorem}{Theorem}
\newtheorem*{thm}{Theorem}
\newtheorem*{lemma}{Lemma}
\newtheorem*{proposition}{Proposition}
\newtheorem*{prop}{Proposition}
\newtheorem*{corollary}{Corollary}
\newtheorem*{cor}{Corollary}
\newtheorem*{claim}{Claim}
\theoremstyle{definition}
\newtheorem*{definition}{Definition}
\newtheorem*{defn}{Definition}
\newtheorem*{remark}{Remark}
\newtheorem*{example}{Example}
\newcommand{\R}{\mathbb R}
\newcommand{\C}{\mathbb C}
\newcommand{\Z}{\mathbb Z}
\newcommand{\Q}{\mathbb Q}
\newcommand{\N}{\mathbb N}
\newcommand{\F}{\mathbb F}
\newcommand{\D}{\mathbb D}
\newcommand{\bB}{\mathbb B}
\newcommand{\bC}{\mathbb C}
\newcommand{\bR}{\mathbb R}
\newcommand{\bZ}{\mathbb Z}
\newcommand{\bN}{\mathbb N}
\newcommand{\bQ}{\mathbb Q}
\newcommand{\bD}{\mathbb D}
\newcommand{\bF}{\mathbb F}
\newcommand{\CP}{\mathbb{CP}}
\newcommand{\RP}{\mathbb{RP}}
\newcommand{\sO}{\mathscr O}
\newcommand{\sI}{\mathscr I}
\newcommand{\sF}{\mathscr F}
\newcommand{\sG}{\mathscr G}
\newcommand{\sH}{\mathscr H}
\newcommand{\sR}{\mathscr R}
\newcommand{\sM}{\mathscr M}
\newcommand{\sV}{\mathscr V}
\newcommand{\sU}{\mathscr U}
\DeclareMathOperator{\Hom}{Hom}
\DeclareMathOperator{\Ext}{Ext}
\DeclareMathOperator{\Tor}{Tor}
\DeclareMathOperator{\Spec}{Spec}
\DeclareMathOperator{\Proj}{Proj}
\DeclareMathOperator{\supp}{supp}
\DeclareMathOperator{\im}{im}
\DeclareMathOperator{\id}{id}
\DeclareMathOperator{\Tr}{Tr}
\DeclareMathOperator{\tr}{tr}
\DeclareMathOperator{\rank}{rank}
\DeclareMathOperator{\ord}{ord}
\DeclareMathOperator{\dist}{dist}
\DeclareMathOperator{\End}{End}
\DeclareMathOperator{\Aut}{Aut}
\DeclareMathOperator{\Gal}{Gal}
\DeclareMathOperator{\vol}{vol}
\DeclareMathOperator{\Cl}{Cl}
\DeclareMathOperator{\coker}{coker}
\DeclareMathOperator{\codim}{codim}
\DeclareMathOperator{\divergence}{div}
\DeclareMathOperator{\Ric}{Ric}
\DeclareMathOperator{\grad}{grad}
\DeclareMathOperator{\Hess}{Hess}
\newcommand{\abs}[1]{\left\lvert#1\right\rvert}
\newcommand{\sabs}[1]{\lvert#1\rvert}
\newcommand{\babs}[1]{\bigl\lvert#1\bigr\rvert}
\newcommand{\Babs}[1]{\Bigl\lvert#1\Bigr\rvert}
\newcommand{\bbabs}[1]{\biggl\lvert#1\biggr\rvert}
\newcommand{\BBabs}[1]{\Biggl\lvert#1\Biggr\rvert}
\newcommand{\norm}[1]{\left\lVert#1\right\rVert}
\newcommand{\snorm}[1]{\lVert#1\rVert}
\newcommand{\bnorm}[1]{\bigl\lVert#1\bigr\rVert}
\newcommand{\Bnorm}[1]{\Bigl\lVert#1\Bigr\rVert}
\newcommand{\bbnorm}[1]{\biggl\lVert#1\biggr\rVert}
\newcommand{\BBnorm}[1]{\Biggl\lVert#1\Biggr\rVert}
\newcommand{\widebar}[1]{\overline{#1}}
\newcommand{\nicefrac}[2]{\frac{#1}{#2}}
\newcommand{\linnprod}[2]{\langle#1,#2\rangle}
\newcommand{\myindex}[1]{#1}
\newcommand{\glsadd}[1]{}
\newcommand{\avoidbreak}{}
\newcommand{\sourceref}[1]{\textnormal{[原文：\nolinkurl{#1}]}}
\newcommand{\thmref}[1]{\sourceref{#1}}
\newcommand{\lemmaref}[1]{\sourceref{#1}}
\newcommand{\propref}[1]{\sourceref{#1}}
\newcommand{\corref}[1]{\sourceref{#1}}
\newcommand{\defnref}[1]{\sourceref{#1}}
\newcommand{\exampleref}[1]{\sourceref{#1}}
\newcommand{\exerciseref}[1]{\sourceref{#1}}
\newcommand{\figureref}[1]{\sourceref{#1}}
\newcommand{\sectionref}[1]{\sourceref{#1}}
\newcommand{\appendixref}[1]{\sourceref{#1}}
\newcommand{\stacksref}[2]{\href{https://stacks.math.columbia.edu/tag/#1}{#2 [#1]}}

\begin{document}
\section*{032\quad 跨紧洞的Hartogs全纯延拓}
\subsection*{来源与计数目标}
Ji\v{r}\'i Lebl, \emph{Tasty Bits of Several Complex Variables}。从作者公开的原生 TeX 正文摘录；获取日期：2026-09-28。使用实际访问的在线正文，不把工作分支冒称冻结的印刷版。
\par\url{https://raw.githubusercontent.com/jirilebl/scv/master/scv.tex}
\par\url{https://www.jirka.org/scv/}
原文位置：\texttt{sec:compactdbar} 中 \texttt{thm:compactdbar} 与 \texttt{sec:hartogsphenom} 的 Hartogs phenomenon 定理及证明。
\par\textbf{本文件只计一个问题：}连通补集情形的 Hartogs 延拓；紧支集 $\bar\partial$ 求解作为同一题的核心支撑，不另计题数。
\subsection*{作者命题与人类证明}

\paragraph{Compactly supported $\bar\partial$-problem (supporting theorem).}
For a smooth function $\psi$, write
\[
\bar\partial\psi=\frac{\partial\psi}{\partial\bar z_1}d\bar z_1+\cdots+\frac{\partial\psi}{\partial\bar z_n}d\bar z_n.
\]
\begin{thm}
Suppose $g$ is a $(0,1)$-form on $\C^n$, $n\geq2$, given by
\[g=g_1d\bar z_1+\cdots+g_nd\bar z_n,\]
where $g_j:\C^n\to\C$ are compactly supported smooth functions satisfying the compatibility conditions
\begin{equation}\label{eq:compatconds}
\frac{\partial g_k}{\partial\bar z_\ell}=\frac{\partial g_\ell}{\partial\bar z_k}\qquad\text{for all }k,\ell=1,2,\ldots,n.
\end{equation}
Then there exists a unique compactly supported smooth function $\psi:\C^n\to\C$ such that $\bar\partial\psi=g$.
\end{thm}
The compactly supported solution is unique: If $\psi_1$ and $\psi_2$ are solutions, then $\bar\partial(\psi_1-\psi_2)=g-g=0$, and so $\psi_1-\psi_2$ is holomorphic. The only holomorphic compactly supported function is $0$, and hence the compactly supported solution $\psi$ is unique.
\begin{proof}
Really there are $n$ smooth functions, $g_1,\ldots,g_n$, so the equation $\bar\partial\psi=g$ consists of the $n$ equations
\[\frac{\partial\psi}{\partial\bar z_k}=g_k,\]
where the functions $g_k$ satisfy the compatibility conditions \eqref{eq:compatconds}. We claim that the following is an explicit solution:
\begin{align*}
\psi(z)&=\frac1{2\pi i}\int_\C\frac{g_1(\zeta,z_2,\ldots,z_n)}{\zeta-z_1}\,d\zeta\wedge d\bar\zeta\\
&=\frac1{2\pi i}\int_\C\frac{g_1(\zeta+z_1,z_2,\ldots,z_n)}{\zeta}\,d\zeta\wedge d\bar\zeta.
\end{align*}
To show that the singularity does not matter for integrability is the same idea as for the generalized Cauchy formula. Let us check that $\psi$ is the solution. We use the generalized Cauchy formula on the $z_1$ variable. Take $R$ large enough so that $g_k(\zeta,z_2,\ldots,z_n)$ is zero when $|\zeta|\geq R$ for all $k$. For every $k$,
\begin{align*}
g_k(z)&=\frac1{2\pi i}\int_{|\zeta|=R}\frac{g_k(\zeta,z_2,\ldots,z_n)}{\zeta-z_1}\,d\zeta\\
&\quad+\frac1{2\pi i}\int_{|\zeta|\leq R}\frac{\frac{\partial g_k}{\partial\bar z_1}(\zeta,z_2,\ldots,z_n)}{\zeta-z_1}\,d\zeta\wedge d\bar\zeta\\
&=\frac1{2\pi i}\int_\C\frac{\frac{\partial g_k}{\partial\bar z_1}(\zeta,z_2,\ldots,z_n)}{\zeta-z_1}\,d\zeta\wedge d\bar\zeta.
\end{align*}
Using the second form of the definition of $\psi$, the compatibility conditions \eqref{eq:compatconds}, and the computation above, we get
\begin{align*}
\frac{\partial\psi}{\partial\bar z_k}(z)
&=\frac1{2\pi i}\int_\C\frac{\frac{\partial g_1}{\partial\bar z_k}(\zeta+z_1,z_2,\ldots,z_n)}{\zeta}\,d\zeta\wedge d\bar\zeta\\
&=\frac1{2\pi i}\int_\C\frac{\frac{\partial g_k}{\partial\bar z_1}(\zeta+z_1,z_2,\ldots,z_n)}{\zeta}\,d\zeta\wedge d\bar\zeta\\
&=\frac1{2\pi i}\int_\C\frac{\frac{\partial g_k}{\partial\bar z_1}(\zeta,z_2,\ldots,z_n)}{\zeta-z_1}\,d\zeta\wedge d\bar\zeta=g_k(z).
\end{align*}
\paragraph{Exercise in the original proof.} Show that we were allowed to differentiate $\psi$ under the integral in the computation above, but only in the second form of the integral.

That $\psi$ has compact support follows because $g_1$ has compact support together with the identity theorem. In particular, $\psi$ is holomorphic for large $z$ since $\bar\partial\psi=g=0$ when $z$ is large. When at least one of $z_2,\ldots,z_n$ is large, then $\psi$ is identically zero simply from its definition.

That is, $\psi$ is zero on the unbounded component of the set where $g=0$, and so $\psi$ has compact support.
\end{proof}

\paragraph{The general Hartogs phenomenon (the counted problem).}
\begin{thm}[Hartogs phenomenon]
Let $U\subset\C^n$ be a domain, $n\geq2$, and let $K\subset\subset U$ be a compact set such that $U\setminus K$ is connected. Every holomorphic $f:U\setminus K\to\C$ extends uniquely to a holomorphic function on $U$.
\end{thm}
The idea of the proof is extending in any way whatsoever and then using the solution to the $\bar\partial$-problem to correct the result to make it holomorphic.
\begin{proof}
First find a smooth function $\varphi$ that is $1$ in a neighborhood of $K$ and is compactly supported in $U$ (exercise below). Let $f_0=(1-\varphi)f$ on $U\setminus K$ and $f_0=0$ on $K$. The function $f_0$ is smooth on $U$ and it is holomorphic and equal to $f$ near the boundary of $U$, where $\varphi$ is $0$. We let $g=\bar\partial f_0$ on $U$, that is, $g_k=\frac{\partial f_0}{\partial\bar z_k}$, and we let $g=0$ outside $U$. As the $g_k$ are identically zero near $\partial U$, we find that each $g_k$ is smooth on $\C^n$. The compatibility conditions \eqref{eq:compatconds} are satisfied because partial derivatives commute. Let us see why $g$ is compactly supported. The only place to check is on $U\setminus K$ as elsewhere we have $g=0$ automatically. Note that $f$ is holomorphic on $U\setminus K$ and compute
\[
\frac{\partial f_0}{\partial\bar z_k}
=\frac{\partial}{\partial\bar z_k}\bigl((1-\varphi)f\bigr)
=(1-\varphi)\frac{\partial f}{\partial\bar z_k}-\frac{\partial\varphi}{\partial\bar z_k}f
=-\frac{\partial\varphi}{\partial\bar z_k}f.
\]
The function $\frac{\partial\varphi}{\partial\bar z_k}$ is compactly supported in $U\setminus K$ by construction. Apply the solution of the compactly supported $\bar\partial$-problem to find a compactly supported function $\psi$ such that $\bar\partial\psi=g$. Set $F=f_0-\psi$. We check that $F$ is the desired extension. First, it is holomorphic:
\[
\frac{\partial F}{\partial\bar z_k}=\frac{\partial f_0}{\partial\bar z_k}-\frac{\partial\psi}{\partial\bar z_k}=g_k-g_k=0.
\]
Next, \exerciseref{exercise:supportofpsi} and the fact that $U\setminus K$ is connected reveal that $\psi$ must be compactly supported in $U$. This means that $F$ agrees with $f$ near the boundary (in particular on an open set) and thus everywhere in $U\setminus K$ since $U\setminus K$ is connected.
\end{proof}

\subsection*{整理说明（非作者正文）}
保留作者的紧支集积分求解、唯一性说明、相容条件和延拓修正全程。标准先修为一元 Cauchy--Pompeiu 公式、可积核下的光滑参数微分、光滑截断函数和恒等定理；原作者把参数微分的合法性及截断函数存在列为练习，本条明示而不伪称附有其证明。末段引用的 \texttt{exercise:supportofpsi} 的无界分支消失性质，在前附作者求解证明的最后一段已经写出；没有把紧支集求解本身当作未展开的先修。

原书的两个示意图 \texttt{fig:dbarcpt-fig}、\texttt{fig:hartogs-fig} 未复制；删去仅指认图中灰白区域的重复说明，保留其前后的恒等定理与无界分支论证。定义 $\bar\partial$ 的公式摘自同节原文；章节提示和计数目标属整理文字。难度依据是将延拓问题转成具紧支集和相容性的非齐次偏微分方程，再构造积分解修正，而不是直接引用 Hartogs 定理。
\subsection*{可修改的简短标注（非作者正文）}
$c$：多变量全纯函数的跨紧洞延拓。

$a$：光滑截断；非齐次 $\bar\partial$ 方程；Cauchy--Pompeiu 积分；相容条件；紧支集；恒等定理。

\texttt{cross\_theory\_status = yes}。把任意光滑延拓的非全纯误差编码为闭的 $(0,1)$-形式，积分求解后消去该误差，利用支集与连通性返回原函数的全纯延拓。位置：前附求解定理及主证明。

全部标注均为非穷尽、可修改初稿；未标注不表示未使用。
\subsection*{许可与修改记录}
原数学正文作者：Ji\v{r}\'i Lebl。作者网站及源书声明允许 CC BY-SA 4.0 或 CC BY-NC-SA 4.0 双许可；本条选择 CC BY-SA 4.0，整理部分亦采用同一许可。\url{https://creativecommons.org/licenses/by-sa/4.0/}。2026-09-28：助手截取题证、调整独立排版，加入来源、范围与初稿标注；不暗示作者认可本次整理。保留的是所用 TeX 正文摘录，未将其称为整书原件。
\end{document}
